\section{The source theorem and its transferable mechanism}\label{sec:src}

\providecommand{\srcrem}{\mathrm{rem}}
\providecommand{\srcB}{\mathcal{B}}
\providecommand{\srcT}{\mathcal{T}}
\providecommand{\srcRR}{\mathcal{R}}
\providecommand{\srctr}{\mathsf{T}}
\begingroup\emergencystretch=3em

This section fixes the statement that is used as input in the rest of the article, describes
how its proof is organised, and isolates those parts of the mechanism that can be stated and
proved without reference to the source. Assertions about the source development describe its
content at the pinned commit; they are quoted, not re-proved. Everything stated as a lemma or
proposition in this section is proved here and does not depend on the input theorem.

\subsection{The input theorem}\label{src:sec:input}
\begingroup\raggedright
\sources{grh_transfer_analysis; grh_transfer_analysis_2026-11-25; grh_lean_transfer_analysis; grh-arithmetic-transfers; grh-lean-transfers; grh-exponent-six-and-transfers; grh_method_transfers_2026-11-25_1; grh_exact_torsion_transfer_2026-11-25; grh_fixed_field_transfer_2026-11-25; grh_structural_transfers_2026-11-25; grh_spanning_and_s5_2026-11-26; grh_arithmetic_transfers_2026-11-26}
\endgroup

The input is the Lean~4 development \repo{} (\code{lsernau/grh-lean}) \cite{Sernau2026GRH} at
the commit
\begin{center}
\code{72fead102f663d5c7a407964a13dc0cf9d977207}.
\end{center}
It is a shorter
rewriting, by L.~Sernau and collaborators, of the argument of L.~Alp\"oge developed with
Claude \cite{Alpoge2026Remark}. The repository contains a readable edition \code{grh} and a
compact edition \code{grh-compact}, both pinned to Lean \code{4.35.0} and to Mathlib revision
\code{56a00585}. The file \code{Statements.lean} is the same
in both editions, imports only Mathlib, and is phrased with Mathlib's Riemann zeta function,
Dirichlet characters and analytically continued Dirichlet $L$-functions. The file
\code{Main.lean} proves five statements: \code{riemannHypothesis} (Mathlib's
\code{RiemannHypothesis}), \code{riemannZeta_eq_zero_re},
\code{LFunction_ne_zero_of_half_lt_re}, \code{LFunction_eq_zero_re} and
\code{LFunction_eq_zero_re_of_isPrimitive}. In conventional notation they say the following.

\begin{theorem}[Input theorem: RH and GRH for Dirichlet $L$-functions]\label{src:thm:dgrh}
\leavevmode
\begin{enumerate}[label=\textup{(\alph*)}]
\item If $\zeta(s)=0$ and $s$ is not a negative even integer, then $\Re s=1/2$. (In the
formal statement the point $s=1$, where the total function has a conventional value, is
excluded separately.)
\item For every integer $N\ge1$, every Dirichlet character $\chi$ modulo $N$ and every
$s\in\C$ with $\Re s>1/2$, one has $L(s,\chi)\ne0$, except that the principal character is
not evaluated at its pole $s=1$.
\item For every integer $N\ge1$ and every Dirichlet character $\chi$ modulo $N$, every zero
of $L(s,\chi)$ in the strip $0<\Re s<1$ has real part $1/2$.
\end{enumerate}
\end{theorem}

The fifth exported statement is the restriction of (c) to primitive characters. The following
points about the scope of \cref{src:thm:dgrh} are used throughout.
\begin{itemize}
\item Imprimitive characters are included in (b) and (c). Their additional Euler factors can
vanish on $\Re s=0$, which is why (c) is stated for the open strip; nothing is asserted for
$\Re s\le0$.
\item The exception at $s=1$ in (b) reflects the pole of the principal $L$-function, which a
total function must be given a conventional value at; it is not an exemption for a zero.
\item Nothing is asserted about the order of a zero on the critical line. A zero at $s=1/2$,
and zeros of any multiplicity on the line, are compatible with the theorem.
\item The functions are Dirichlet $L$-functions over $\Q$. The statement does not quantify
over Hecke characters of other number fields, Artin representations of dimension greater
than one, Dedekind zeta functions of nonabelian fields, or automorphic $L$-functions. The
proof uses CM fields, unitary groups, Hecke values and Galois representations, but this does
not enlarge the quantifiers of its conclusion.
\end{itemize}
Throughout the article \emph{Dirichlet GRH} means \cref{src:thm:dgrh}.

The repository reports that the five theorems depend only on the axioms \code{propext},
\code{Classical.choice} and \code{Quot.sound}, and independent builds and kernel replays of
both editions at this commit are reported by Sutherland, Rubinstein and Sonoda
\cite{Sutherland2026Verification,Rubinstein2026Check,Sonoda2026Build}. The present article
takes \cref{src:thm:dgrh} as input and does not re-verify it: none of these builds or checks
was repeated. A result marked ``\stG'' is a theorem exactly to the extent that
\cref{src:thm:dgrh} is; a result marked ``\stU'' does not use it.

\subsection{Architecture of the source proof}\label{src:sec:mechanism}
\begingroup\raggedright
\sources{grh_transfer_analysis; grh_transfer_analysis_2026-11-25; grh_lean_transfer_analysis; grh_beyond_dirichlet_grh_2026-11-25; grh_beyond_dirichlet_2026-11-25; grh_fixed_field_transfer_2026-11-25; grh_exact_torsion_transfer_2026-11-25; grh_method_transfers_2026-11-25; grh_method_transfers_2026-11-25_1; grh_method_transfers_2026-11-25_2; grh_method_transfers_orthogonal_jets_2026-11-26; grh_sharp_transfers_2026-11-25; grh_structural_transfers_2026-11-25; grh_transfer_bridges_2026-11-25; grh_transfer_extensions_2026-11-26; grh_transfer_synthesis_2026-11-25; grh_transfer_advances_2026-11-26; grh_arithmetic_transfers_2026-11-26; grh_nonabelian_transfer_2026-11-25; grh_spanning_and_s5_2026-11-26; grh_ideal_densities_and_transfer_conditions_2026-11-26; grh_a6_transfer_2026-11-25; grh_cyclic_generator_transfer_2026-11-25; grh_deficit_transfer_2026-11-25; grh_kummer_window_transfer_2026-11-25; grh_relative_prime_races_2026-11-25; grh_sylow_transfer_2026-11-25; psl2_grh_transfer; grh-arithmetic-transfers; grh-lean-transfers; grh_new_results_2026-11-26; symplectic_rank_blindness_2026-11-26; squarefree_prime_products}
\endgroup

The proof has an arithmetic part, which bounds a smoothed prime sum from above, and an
analytic part, which turns that one-sided bound into the absence of zeros. Declarations are
quoted from the readable edition unless the compact edition is named.

\paragraph{The prime statistic.}
Let $q\ge1$ be a modulus and $d$ an odd auxiliary prime with quadratic character
$\chi_d(n)=(n/d)$. Put
\begin{equation}\label{src:eq:R}
\lambda_{d,q}(n)=\Lambda(n)\,\mathbf 1_{n\equiv1\ (\mathrm{mod}\ q)}\bigl(1-\chi_d(n)\bigr),\qquad
R_{d,q}(x)=\sum_{n\le x}\lambda_{d,q}(n)\Bigl(\frac1n-\frac1x\Bigr)-\frac{\log x}{\varphi(q)}-c_{d,q},
\end{equation}
where $\Lambda$ is the von Mangoldt function and $c_{d,q}$ is the constant for which
$R_{d,q}(x)\to0$. The coefficients $\lambda_{d,q}(n)$ are nonnegative and supported on prime
powers $n\equiv1\pmod q$ with $\chi_d(n)\ne1$; in the admissible setting $\chi_d$ is the
quadratic character of the field $K=\Q(\sqrt{-d})$. The ramp kernel
$1/n-1/x$ vanishes at $n=x$ and puts the factor $1/(s(s-1))$ into the Mellin transform. In
the arithmetic construction the modulus is a multiple $q_0$ of $18480$ and $d$ is a prime
with $d\equiv q_0/2-1\pmod{q_0}$, so that $d\equiv7\pmod 8$; the auxiliary field is the CM
field $E=K(\zeta_{q_1})$, which is abelian over $\Q$. These restrictions are internal. The
final step raises an arbitrary modulus $N$ to $18480N$ and uses change of level, the extra
Euler factors being harmless in the half-plane concerned.

\paragraph{Three estimates for one lattice.}
The construction uses definite unitary groups over $E$, derivatives of CM theta functions and
chains of $m$ successive theta lifts between spaces of algebraic automorphic forms carrying
integral structures. Let $D$ be the dimension of the target space, $n$ the size of a linearly
independent family selected from images of integral source forms, $V$ an upper bound for the
logarithm of the Gram determinant (squared covolume) of that family, and $\delta$ a bound for
the logarithmic index needed to make the family integral at all finite places. Write $M$ for
the common main amount (\code{mainAmount}), $\srcrem$ for the unnormalised prime remainder
(\code{totalRemainder}), $f$ for the number of infinite places of $E$, $a=xb$ for the order
of derivative and $b$ for a much larger auxiliary weight. The reading guide of the source
names three inputs.
\begin{enumerate}[label=\textup{(\Alph*)}]
\item \emph{Short forms.} High derivatives of theta functions are highly divisible at
selected primes, and the divisibility records the weighted prime sum. After normalisation
this gives integral source forms whose images are short, with controlled index:
\[
V\le n\bigl(M-m\,\srcrem+L\bigr),\qquad \delta\le D\,I,\qquad
L=fa\log^4b,\quad I=10\log(1+2b^2)\,fa\log^3b .
\]
\item \emph{Enough forms.} A second-moment argument shows that a controlled number of
translates of the images span almost all of the target: $(1-4/b^2)D\le n\le D$, and the number
of translates times the source dimension is at most $5D\log(1+2b^2)$.
\item \emph{Integral forms are long.} A product-formula argument with determinants of chains
of lifts, twists and denominator bounds gives a lower bound for the same corrected covolume:
$D(M-3L)\le V+2\delta$.
\end{enumerate}
The three statements refer to the same lattice of target forms with the same integral
structures, and that is what makes them comparable.

Input (C) is an instance of an abstract statement about twisted lattices, the theorem
\code{covolume_add_index_nonneg} of the namespace \code{New.TwistedLattice}. It concerns
linear maps $\Theta_\xi$,
indexed by twists $\xi$, from one complex vector space to inner-product spaces, integral
lattices at the $p$-adic identifications, and rational positive-definite Gram data at the base
twist. If $n$ chosen vectors are completed to a basis by $n'$ further vectors, it gives
\begin{equation}\label{src:eq:twisted}
0\le c+2\delta+n'\bigl(b_{\mathrm{short}}+b_{\mathrm{denom}}\bigr)+b_{\mathrm{volume}},
\end{equation}
where $c$ bounds the logarithmic Gram determinant of the chosen vectors, $\delta$ their index,
$b_{\mathrm{short}}$ and $b_{\mathrm{denom}}$ the squared norms and denominators of the
completing vectors, and $b_{\mathrm{volume}}$ the heights of ratios of lift determinants.
Writing $Z(\xi)$ for a lift determinant and $\Gamma(\xi)$ for the Gram determinant of the
lifted basis, the proof factors $\Gamma(1)$ at each allowed local twist $\xi$ into an integral
term, a denominator term and $Z(1)/Z(\xi)$; the height bounds add, and archimedean norm bounds
together with nonvanishing of the full determinant give \eqref{src:eq:twisted}. The favourable
twist may be chosen separately at each $p$-adic embedding, while the archimedean size is read
on the untwisted map; \cref{src:sec:quantifiers} records the compatibility conditions this
imposes. The existence of the required twists is where the congruence and Selmer machinery
of the source enters: an anticyclotomic twisting theorem
(\code{New.Selmer.exists_anticyclotomic_twist}) built on a selection lemma for $p$-power
roots of unity (\code{RootsOfUnity.lean}), developed in \cref{sec:tors,sec:cost,sec:sel}.

\paragraph{Cancellation of the main term.}
Combining (A)--(C) and using $M\ge0$ and $n\le D$ gives $nm\,\srcrem\le D(4L+2I)$; since
$n\ge D/2$ and the scales are widely separated, the source concludes
(\code{New.star_inequality})
\begin{equation}\label{src:eq:star}
m\,\srcrem\le fa(\log b)^5 .
\end{equation}
This step is elementary once (A)--(C) are available, and its general form is
\cref{src:prop:amplification}. The compact edition writes the same comparison with budget
$\srcB=fa\log^3b$, loss $L=\srcB\log b$ and index bound $\delta\le5D\log(1+b^2)\srcB$, replaces
the sign condition $M\ge0$ by a one-sided bound on the top height, and proves
$m\,\srcrem\le(\log b)^2\srcB$; this is the case of \cref{src:prop:top} recorded after it.

\paragraph{From the main inequality to the prime sum.}
An exact finite identity expresses the remainder through \eqref{src:eq:R}: in the compact
edition (\code{Data.remainder_eq}),
$\srcrem=a\bigl(\tau(R_{d,q_0}(x)-R_{d,q_0}(x_1))+\tau_0R_{d,1}(x_1)\bigr)$, where $\tau$ and
$\tau_0$ count the thick and the top places and $x_1$ is a second cutoff. The identity rests
on a count of places in cyclotomic classes. From \eqref{src:eq:star} the readable edition
obtains
\begin{equation}\label{src:eq:cutoffs}
R_{d,q_0}(x)\le\frac{\varphi(q_0)(\log b)^6}{m}+\abs{R_{d,q_0}(x_1)}+(\log b)\abs{R_{d,1}(x_1)} .
\end{equation}
The parameters satisfy $\sqrt x\le m\le2\sqrt x$, $\log b\le2x^{\eps}$ and
$\log x_1\ge x^{\eps/12}$. The two terms at $x_1$ are bounded by the prime number theorem for
fixed characters; no GRH input is used there. For integers $x$ and $0<\eps<1/100$ this gives
\begin{equation}\label{src:eq:onesided-int}
R_{d,q_0}(x)\le(64\varphi(q_0)+3)\,x^{-1/2+6\eps}
\end{equation}
eventually, and a comparison of a real $x$ with its integer part gives, for every $\eps>0$,
\begin{equation}\label{src:eq:onesided}
R_{d,q_0}(x)\le x^{-1/2+\eps}\qquad\text{for all sufficiently large real }x .
\end{equation}
The threshold depends on $d$ and $q_0$; no uniform onset is asserted. The compact edition
chooses the second scale differently: a window lemma
(\code{IsPolar.exists_window_sub_mul_le}) supplies, for all large $t$ and every coefficient
$0\le X\le e^t$, a point $v\in[t,t^2]$ with $E_0(v)-XE_1(v)\le e^{-t}$, where $E_0$ and $E_1$
are the two errors $v\mapsto R_{d,q_0}(e^v)$ and $v\mapsto R_{d,1}(e^v)$; the resulting
exponent is $-1/2+8\eps$. Sharper windows are given in \cref{ana:thm:window}.

The exponent $1/2$ comes from the length $m\asymp\sqrt x$ of the chain of lifts. In the
announcement of the original proof each additional lift is compared with additional
multiplicity in Stepanov's method \cite{Alpoge2026Remark}. The bound is one-sided: an upper
bound for $R_{d,q_0}$, not an estimate for its absolute value.

\paragraph{From the one-sided bound to zeros.}
Put $E(v)=R_{d,q}(e^v)$ and $\srcT(s)=\int_0^\infty E(v)e^{(1-s)v}\,\dd v$. The transform has
a meromorphic continuation to $\Re s>0$ with at most simple poles
(\code{isPolar_primeRemainder}). Writing $L^\circ(s,\chi)$ for the Dirichlet $L$-function
with its principal pole removed, the residue at a point $\rho$ other than the separately
treated ones is
\begin{equation}\label{src:eq:residue}
\Res_{s=\rho}\srcT(s)=
\frac{\sum_{\chi\bmod q}\ord_\rho L^\circ(s,\chi)-\sum_{\chi\bmod q}\ord_\rho L^\circ(s,\chi\chi_d)}
{\varphi(q)\,\rho(1-\rho)} .
\end{equation}
The numerator is a difference, so nonnegative prime coefficients alone do not prevent a zero
of an untwisted function from being cancelled by a zero of a twisted one. The source removes
the second sum by a choice of $d$. Given a hypothetical zero $\rho$ with $\Re\rho>\beta>1/2$,
put $T=\abs{\Im\rho}$ and $\delta_\beta=(\beta-1/2)/9$. A zero-density estimate
(\code{zero_density_quadTwist}) bounds by $O_{q,\beta,T}(X^{1-\delta_\beta/2})$ the number of
permitted primes $d\le X$ for which some $L(s,\chi\chi_d)$, $\chi$ modulo $q$, has a zero in
$\Re s\ge\beta$, $\abs{\Im s}\le T$; hence some permitted $d$ has no such zero
(\code{exists_entireL_quadTwist_ne_zero}). For that $d$ the surviving multiplicities in
\eqref{src:eq:residue} are nonnegative integers in the box. At real $s\in(\beta,1)$ the
denominator $s(1-s)$ is positive, and for $s>1$ the functions $L^\circ$ do not vanish, so
every real residue to the right of $\beta$ is nonnegative. Landau's theorem on Laplace
transforms of nonnegative functions \cite[Lemma~15.1]{MontgomeryVaughan2007Multiplicative},
in the form \code{New.IsPolar.landau}, then shows that the eventual bound
$E(v)\le e^{(\beta-1)v}$ excludes every pole of $\srcT$ with real part greater than $\beta$.
The residue at $\rho$ is therefore zero, its numerator is a sum of nonnegative integers
containing the multiplicity of the hypothetical zero, and this is a contradiction. This gives
(b) of \cref{src:thm:dgrh} for the admissible moduli; the functional equation gives (c). The
order of quantifiers is: first $\rho$, then $d$, then the eventual estimate for that $d$. The
criterion (\code{New.LFunction_ne_zero_of_primeRemainder_le}) accordingly assumes
\eqref{src:eq:onesided} for all auxiliary primes in the prescribed progression, and a bound
for one convenient twist does not replace this. The one-sided step and all its retained
variants, including the weakening to a sufficiently large set of auxiliary primes, are the
subject of \cref{sec:ana}; see \cref{ana:thm:onesided}. Related one-sided criteria predate the
source; the compact edition points to the treatment of Suzuki \cite{Suzuki2025Variants}.

\paragraph{Declarations.}
The stages correspond to the following declarations, quoted elsewhere in the article.
\begin{center}\footnotesize
\begin{tabular}{@{}>{\raggedright\arraybackslash}p{0.17\linewidth}>{\raggedright\arraybackslash}p{0.41\linewidth}>{\raggedright\arraybackslash}p{0.36\linewidth}@{}}
\toprule
Stage & Readable edition & Compact edition\\
\midrule
Short forms, index & \code{span_short_integralRamified_eq_top}, \newline\code{indexLe_normalizedPoincare} & \code{Data.short_forms}\\
Enough forms & \code{enough_forms} & \code{Data.enough_forms}\\
Long forms & \code{integral_forms_long} & \code{Data.integral_forms_long}\\
Twisted lattice & \code{covolume_add_index_nonneg}, \code{HeightLe} & \\
Fitted twists & \code{FitTwistsExist.lean} & \code{exists_isFitFamily}\\
Cancellation & \code{star_inequality} & \code{star_inequality}\\
Prime sum & \code{primeRemainder_le}, \newline\code{eventually_primeRemainder_natCast_le} & \code{Data.remainder_eq}, \newline\code{eventually_primeError_le_rpow}\\
Transform & \code{isPolar_primeRemainder} & \code{isPolar_primeError}\\
Twist separation & \code{zero_density_quadTwist}, \newline\code{exists_entireL_quadTwist_ne_zero} & \code{exists_entireL_quadTwist_ne_zero}\\
Landau & \code{IsPolar.landau} & \code{IsPolar.landau}\\
Criterion & \code{LFunction_ne_zero_of_primeRemainder_le} & \code{entireL_ne_zero_of_primeError_le}\\
\bottomrule
\end{tabular}
\end{center}

\paragraph{Ingredients used on their own.}
Several internal results of the source do not depend on its final theorem and are developed
separately in this article: the selection of $p$-power roots of unity for finitely many
integral polynomials (\code{RootsOfUnity.lean}; \cref{sec:tors,sec:cost}); the anticyclotomic
twisting theorem and the bound \code{New.exists_length_shaOdd_le} for odd cocycle modules
(\cref{sec:sel}); the second-moment spanning lemmas and the positive-type matrix lemma
\code{New.card_mul_sum_det_eq_zero_le} (\cref{sec:spn}); the weighted Behrend-type lemma
\code{exists_three_term_free} (\cref{sec:gpp,sec:gpi}); the sharpness lemma for one curve,
\code{New.exists_monomial_max_norm_coeff} (\cref{sec:lub}); and the estimates for principal
submatrices and Schur complements in the file \code{HermitianMatrices.lean}
(\cref{sec:gra}).

\subsection{The determinant comparison}\label{src:sec:comparison}
\begingroup\raggedright
\sources{grh_transfer_analysis; grh_transfer_analysis_2026-11-25; grh_lean_transfer_analysis; grh_method_transfers_orthogonal_jets_2026-11-26; grh_transfer_advances_2026-11-26; grh_transfer_synthesis_2026-11-25; grh_method_transfers_2026-11-25; grh_a6_transfer_2026-11-25; grh_relative_prime_races_2026-11-25; grh_fixed_field_transfer_2026-11-25; grh_transfer_bridges_2026-11-25}
\endgroup

The cancellation step can be separated completely from the arithmetic. In the following
statements $D$ plays the role of the target dimension, $n$ of the rank of the selected family,
$m$ of the number of lifts, $M$ of the common main term, $R$ of the remainder to be bounded,
$V$ of a logarithmic squared covolume and $\delta$ of a logarithmic index; $U$, $L$ and $I$
are the losses in the upper bound, in the lower bound and in the index.

\begin{proposition}[Determinant amplification; \stU]\label{src:prop:amplification}
Let $D>0$, $m>0$, $0<\theta\le1$ and $\theta D\le n\le D$. Let $L,I\ge0$ and let
$M,U,R,V,\delta$ be real numbers with
\begin{equation}\label{src:eq:sandwich}
D(M-L)\le V+2\delta,\qquad V\le n(M-mR+U),\qquad \delta\le DI .
\end{equation}
Then
\begin{equation}\label{src:eq:amp-exact}
nmR\le nU+D(L+2I)-(D-n)M .
\end{equation}
Moreover:
\begin{enumerate}[label=\textup{(\roman*)}]
\item if $M\ge0$, then $R\le\dfrac{U+(L+2I)/\theta}{m}$;
\item if $U\ge0$ and $M+U\ge0$, then $R\le\dfrac{U+L+2I}{\theta m}$.
\end{enumerate}
\end{proposition}

\begin{proof}
The three hypotheses give $D(M-L)\le n(M-mR+U)+2DI$, which rearranges to
\eqref{src:eq:amp-exact}. For (i), the term $-(D-n)M$ is nonpositive because $n\le D$ and
$M\ge0$; dividing $nmR\le nU+D(L+2I)$ by $nm>0$ and using $D/n\le1/\theta$ and $L+2I\ge0$ gives
the claim. For (ii), rewrite the right-hand side of \eqref{src:eq:amp-exact} as
$(n-D)(M+U)+D(U+L+2I)$. The first term is nonpositive, so $nmR\le D(U+L+2I)$; divide by $nm$
and use $n\ge\theta D$ and $U+L+2I\ge0$.
\end{proof}

In (i) the loss $U$ may be negative, and $\delta$ may be negative in both parts. A Lean
statement of (i), \code{GrhTransfer.sharp_determinant_sandwich}, is delivered in
\artifact{grh_transfer_analysis/lean/Transfer.lean}. With the constants of the
readable edition, namely $U=L_0$ and $L=3L_0$ with $L_0=fa\log^4b$ the loss of (A) and (C),
\eqref{src:eq:amp-exact} together with $M\ge0$ and $n\le D$ gives the inequality
$nm\,\srcrem\le D(4L_0+2I)$ used for \eqref{src:eq:star}.

The compact edition does not use a sign condition on the main term. Its comparison is the
following.

\begin{proposition}[Comparison with a one-sided top height; \stU]\label{src:prop:top}
Let $d>0$, $0<n\le d$ and $d-n\le\kappa d$ with $\kappa\ge0$. Let $J,T\ge0$ and let
$h,L,m,R,V,\delta$ be real numbers. Put $h_{\mathrm{top}}=h-mR+L$ and suppose
\[
d(h-3L)\le V+2\delta,\qquad V\le n\,h_{\mathrm{top}},\qquad \delta\le dJ,\qquad
-h_{\mathrm{top}}\le T .
\]
Then $mR\le4L+\kappa T+2J$.
\end{proposition}

\begin{proof}
Write $n\,h_{\mathrm{top}}=d\,h_{\mathrm{top}}-(d-n)h_{\mathrm{top}}\le d\,h_{\mathrm{top}}+(d-n)T
\le d\,h_{\mathrm{top}}+\kappa dT$. Hence
$d(h-3L)\le d(h-mR+L)+\kappa dT+2dJ$. Divide by $d$ and cancel $h$.
\end{proof}

No bound for $\abs{h_{\mathrm{top}}}$ is needed, only the one-sided one. In the compact
edition $\kappa=4/b^2$, $T=b^2x^{99}$, $J=5\log(1+b^2)\srcB$ and $L=\srcB\log b$ with
$\srcB=fa\log^3b$, so that
\[
mR\le4L+4x^{99}+10\log(1+b^2)\srcB\le(\log b)^2\srcB=fa(\log b)^5 .
\]

\begin{corollary}[Completion cost; \stU]\label{src:cor:completion}
In \cref{src:prop:amplification} assume $M\ge0$ and replace the first inequality of
\eqref{src:eq:sandwich} by
\[
D(M-L_0)-(D-n)K\le V+2\delta
\]
with $L_0,K\ge0$, and assume $n\ge(1-\eta)D$ with $0\le\eta<1$. Then
\[
R\le\frac{U+(L_0+2I+\eta K)/(1-\eta)}{m} .
\]
\end{corollary}

\begin{proof}
Since $D-n\le\eta D$, the hypotheses of \cref{src:prop:amplification} hold with
$L=L_0+\eta K$ and $\theta=1-\eta$. Apply (i).
\end{proof}

A spanning fraction is therefore useful only together with an affordable cost $K$ for each
missing dimension. The next statement tracks this cost in the form in which it arises from
\eqref{src:eq:twisted}, with $r=n'$, $K=b_{\mathrm{short}}+b_{\mathrm{denom}}$ and
$b_{\mathrm{volume}}\le-DA+B$.

\begin{proposition}[Rank deficit versus completion cost; \stU]\label{src:prop:rank}
Let $D=n+r$ with $n>0$ and $r\ge0$, let $m>0$, and let $c,\delta,A,R,U,K,B$ be real numbers
with
\[
c\le n(A-mR+U),\qquad 0\le c+2\delta+rK-DA+B .
\]
Then
\begin{equation}\label{src:eq:rank}
nmR\le nU+2\delta+B+r(K-A) .
\end{equation}
If $A\ge0$, then
\[
R\le\frac{U+(2\delta+B)/n+(r/n)K}{m} .
\]
\end{proposition}

\begin{proof}
Substitute the bound for $c$ into the second inequality, expand $D=n+r$ and cancel $nA$; this
is \eqref{src:eq:rank}, and no sign assumption has been used. If $A\ge0$, discard $-rA\le0$ and
divide by $nm$.
\end{proof}

In \eqref{src:eq:rank} the deficit $r$ is charged by $K-A$, the excess of the completion cost
over the main amount, not by $K$ alone. An extension of the method must therefore achieve a
small rank deficit, or cheaper completing vectors, or a compensating estimate for the term
$-rA$; \cref{sec:spn} makes the corresponding weakening of the source hypothesis explicit.

\begin{remark}[Rates]\label{src:rem:rates}
Suppose a construction provides, at scale $x$, an amplification $m\ge x^{\alpha}$, a spanning
fraction $\theta\ge x^{-\lambda}$ and a normalised loss $U+L+2I\le x^{u+o(1)}$, any completion
cost being included in the losses. Then \cref{src:prop:amplification} gives
$R(x)\le x^{-\alpha+\lambda+u+o(1)}$, and square-root strength requires
$\alpha-\lambda-u\ge1/2$. With the twist-selection step of the source this excludes zeros
only to the right of $\max\{1/2,\,1-\alpha+\lambda+u\}$: the floor $1/2$ belongs to that step
and is not lowered by increasing $m$.
\end{remark}

\begin{remark}[Averaging]\label{src:rem:averaging}
With $\alpha=n/D$, the bound in \cref{src:prop:amplification}(i) involves the combined
quantity $(L+2I)/\alpha$. If the comparison is available only on average over a family of
constructions, separate bounds for the average of $\alpha$ and for the average of $L+2I$ do
not control the average of $(L+2I)/\alpha$; an averaged version must retain the combined loss.
The averaged analytic criteria of \cref{sec:ana} are formulated accordingly.
\end{remark}

\subsection{Integral Gram determinants with vanishing minors}\label{src:sec:gram}
\begingroup\raggedright
\sources{grh_new_results_2026-11-26; gram-euler-and-zero-filters}
\endgroup

Input (C) is a lower bound for the Gram determinant of an integral family, and inputs (A) and
(B) ask for many short independent forms. A finite-dimensional question isolates the tension
between these requirements: how small can the Gram determinant of an integer matrix be when
many of its maximal minors are nonzero? Cauchy--Binet and integrality give one lower bound; the
Pl\"ucker relations between the minors force a stronger one, with an explicit penalty for
vanishing minors.

\begin{theorem}[Arithmetic determinant amplification; \stU]\label{src:thm:gram}
Let $A\in\Z^{r\times n}$ have rank $r$, where $n\ge5$ and $2\le r\le n-2$. For an $r$-element
set $B\subseteq\set{1,\dots,n}$ let $A_B$ be the square matrix of the columns in $B$, in
increasing order, and put
\[
 N=\binom nr,\qquad z=\#\set{B:\det A_B=0},\qquad
 D=\gcd\set{\abs{\det A_B}:\det A_B\ne0}.
\]
Then
\begin{equation}\label{src:eq:gram}
 \frac{\det(AA^{\srctr})}{D^2}\ \ge\ \max\Bigl\{N-z,\ \frac85N-4z,\ \frac{12}5N-12z\Bigr\}.
\end{equation}
In particular, if every $r$ columns are linearly independent,
$\det(AA^{\srctr})\ge\tfrac{12}5\binom nrD^2$. The coefficient $12/5$ cannot be increased
uniformly over this class, and each of the three affine bounds in \eqref{src:eq:gram} is
attained.
\end{theorem}

The theorem does not assert that the maximum in \eqref{src:eq:gram} is optimal for every fixed
triple $(r,n,z)$. Cauchy--Binet and integrality alone give only $N-z$. On four labels the
relevant relation is
\begin{equation}\label{src:eq:plucker}
 p_{12}p_{34}-p_{13}p_{24}+p_{14}p_{23}=0,
\end{equation}
which holds for the two-row minors of any matrix by direct expansion.

\begin{lemma}[Local Pl\"ucker bounds; \stU]\label{src:lem:plucker}
\begin{enumerate}[label=\textup{(\roman*)}]
\item If the six integers in \eqref{src:eq:plucker} are nonzero, their squares sum to at least
$9$.
\item If ten nonzero integers $p_{ij}$, $1\le i<j\le5$, satisfy \eqref{src:eq:plucker} on every
four-element set of labels, their squares sum to at least $24$.
\item For ten integers $p_{ij}$, $1\le i<j\le5$, satisfying \eqref{src:eq:plucker} on every
four labels, let $w$ be the number of zeros and $V$ the sum of squares. Then
$V\ge\max\set{10-w,\ 16-4w,\ 24-12w}$.
\end{enumerate}
\end{lemma}

\begin{proof}
(i) If all six absolute values were one, the left side of \eqref{src:eq:plucker} would be a
signed sum of three ones and could not vanish; so one square is at least four.

(ii) Suppose the sum is at most $23$. No coordinate has absolute value four or more, since
that contributes at least $16+9=25$, and at most one has absolute value three, since two
contribute at least $18+8=26$. If the edge $\set{a,b}$ has absolute value three, at most one
edge has absolute value two, because relative to the baseline $10$ the three adds $8$ and two
twos would add $6$. Choose four labels containing $\set{a,b}$ and avoiding that edge, by
omitting one of its endpoints outside $\set{a,b}$; the three absolute products in the relation
are then $3,1,1$, which have no vanishing signed sum. So all absolute values are one or two.
Call an edge heavy if its absolute value is two. On any four labels the absolute products lie
in $\set{1,2,4}$, and a signed sum of three nonzero reals vanishes only if one absolute value is
the sum of the other two, so the product multiset is $\set{2,1,1}$ or $\set{4,2,2}$, with one or
four heavy edges among the six. There are five four-label subsets and each edge lies in three of
them. If $h$ edges are heavy and $k$ subsets have four heavy edges, double counting gives
$3h=(5-k)+4k=5+3k$, which is impossible modulo three.

(iii) Integrality gives $V\ge10-w$, and (ii) gives $V\ge24$ for $w=0$. If $w=1$, delete an
endpoint of the zero edge: the six coordinates on the remaining four labels are nonzero and
contribute at least $9$ by (i), and the three other nonzero coordinates at least $3$, so
$V\ge12$. These cases prove $V\ge24-12w$, whose right side is nonpositive for $w\ge2$, and
$V\ge16-4w$ for $w\le1$; for $w\ge2$, $10-w\ge16-4w$.
\end{proof}

\begin{proof}[Proof of \cref{src:thm:gram}]
Put $p_B=\det(A_B)/D\in\Z$. By Cauchy--Binet,
$S:=\det(AA^{\srctr})/D^2=\sum_{\abs B=r}p_B^2$. First let $r\le n-3$. For disjoint sets
$I,J$ with $\abs I=r-2$ and $\abs J=5$, and distinct $a,b\in J$, let
$q_{ab}=D^{-1}\det(A_{I,a,b})$, the columns of $I$ in increasing order followed by $a,b$. These
are signed versions of $p_{I\cup\set{a,b}}$ and satisfy the four-label relations: if the
columns of $I$ are dependent all $q_{ab}$ vanish, and otherwise an invertible real row
operation sends them to the first $r-2$ standard basis vectors, and expansion along those
columns writes the $q_{ab}$ as a common real multiple of two-row minors. Let $V_{I,J}$ and
$w_{I,J}$ be the local sum of squares and number of zeros. There are
$\binom n{r-2}\binom{n-r+2}5$ local systems, and each $p_B$ occurs in exactly
$t=\binom r2\binom{n-r}3>0$ of them: choose the two elements of $B$ outside $I$, then the three
elements of $J$ outside $B$. Hence $10\binom n{r-2}\binom{n-r+2}5=Nt$,
$\sum_{I,J}V_{I,J}=tS$ and $\sum_{I,J}w_{I,J}=tz$. Summing each of the three affine
inequalities of \cref{src:lem:plucker}(iii) and dividing by $t$ gives the three bounds.

If $r=n-2$, the complementary coordinates $p^*_{ij}=\pm p_{\set{1,\dots,n}\setminus\set{i,j}}$,
with the signs of the complementary permutation, are, up to one common nonzero real factor, the
two-row minors of a basis of the two-dimensional orthogonal complement of the row space of $A$,
by Laplace expansion. So they are integers satisfying the rank-two relations, with the same
sum of squares and number of zeros, and $\binom n2=N$; apply the case $r=2$. The examples below
give the sharpness assertions.
\end{proof}

\begin{example}[Equality cases]\label{src:ex:gram}
Let $e_1=(1,0)^{\srctr}$, $e_2=(0,1)^{\srctr}$, $u=(1,1)^{\srctr}$, $v=(1,-1)^{\srctr}$
and $w=(1,2)^{\srctr}$. The following matrices have $r=2$, $n=5$, $N=10$ and $D=1$.
\begin{center}\small
\begin{tabular}{cllc}
\toprule
$z$ & columns of $A$ & $AA^{\srctr}$ & $\det(AA^{\srctr})$\\
\midrule
$0$ & $e_1,e_2,u,v,w$ & $\left(\begin{smallmatrix}4&2\\2&7\end{smallmatrix}\right)$ & $24$\\
$1$ & $e_1,e_1,e_2,u,v$ & $\left(\begin{smallmatrix}4&0\\0&3\end{smallmatrix}\right)$ & $12$\\
$2$ & $e_1,e_1,e_2,e_2,u$ & $\left(\begin{smallmatrix}3&1\\1&3\end{smallmatrix}\right)$ & $8$\\
$3$ & $e_1,e_1,e_1,e_2,u$ & $\left(\begin{smallmatrix}4&1\\1&2\end{smallmatrix}\right)$ & $7$\\
\bottomrule
\end{tabular}
\end{center}
The first attains $12N/5-12z$, the second both nontrivial bounds, the third $8N/5-4z$ and the
baseline, and the fourth the baseline alone. The rows of
\[
 \begin{pmatrix}-1&-1&1&0&0\\-1&1&0&1&0\\-1&-2&0&0&1\end{pmatrix}
\]
form an integral basis of the kernel of the first example; this rank-three matrix has ten
nonzero maximal minors with gcd one and Gram determinant $24$. Exact integer arithmetic confirms
all five values, Cauchy--Binet and every maximal minor
(\artifact{grh_new_results_2026-11-26/artifacts/determinant_results.json}).
\end{example}

With $\delta=z/N$ the normalized bound is $12/5-12\delta$ for $\delta\le1/10$,
$8/5-4\delta$ for $1/10\le\delta\le1/5$, and $1-\delta$ for $\delta\ge1/5$.

\begin{corollary}[Short integral columns are limited; \stU]\label{src:cor:gram}
Under the hypotheses of \cref{src:thm:gram} let $E=\sum_{j=1}^n\norm{A_j}^2$ be the total
squared length of the columns and $F$ the right side of \eqref{src:eq:gram}. Then
$E\ge rD^{2/r}F^{1/r}$. If every maximal minor is nonzero and every column has length at most
$H$, then $(nH^2/r)^r\ge\tfrac{12}5\binom nrD^2$.
\end{corollary}

\begin{proof}
The eigenvalues of $AA^{\srctr}$ are positive with sum $E$, so the arithmetic--geometric mean
inequality gives $(E/r)^r\ge\det(AA^{\srctr})$; apply \cref{src:thm:gram}, and $E\le nH^2$.
\end{proof}

A common determinantal divisor, many independent column subsets and a small total squared
length are therefore competing requirements. This does not by itself improve the comparison of
\cref{src:sec:comparison}: an application would have to match the forms and lattices of the
source to the hypotheses and the normalization used here. The weaker fact that a generic
two-row integer matrix with five columns has a minor of absolute value at least three is known
from the study of generic integer matrices with bounded maximal minors
\cite{KriepkeKyureghyanSchymura2023Size,KriepkeSchymura2025Generic}; the aggregate
sum-of-squares inequality and its dependence on vanishing minors are the content of
\cref{src:thm:gram}.

\paragraph{The rank-two minimum for many columns.} For $r=2$ and $z=0$ the least value is known up to
$O(n^2)$. For $A=[v_1\ \cdots\ v_n]\in\Z^{2\times n}$ with $\det(v_i,v_j)\ne0$ for $i\ne j$ put
\[
 D(A)=\gcd_{i<j}\abs{\det(v_i,v_j)},\qquad E(n)=\min_A\frac{\det(AA^{\srctr})}{D(A)^2},
\]
so that \cref{src:thm:gram} gives $E(n)\ge\frac{12}5\binom n2$ for $n\ge5$, with equality at $n=5$. Let
$w_j=\varphi(j)/j$ and
\begin{equation}\label{src:eq:gramPhi}
 \Phi(T)=\frac4{3T}\sum_{1\le j\le T}w_j\Bigl(1-\frac{j^2}{T^2}\Bigr)^{3/2}\qquad(T>0),
\end{equation}
the empty sum being zero, and $\Phi_*=\max_{T>0}\Phi(T)$.

\begin{theorem}[Asymptotic rank-two minimum; \stU]\label{src:thm:gramasym}
The maximum $\Phi_*$ is attained at a unique $T_*\in(6,7)$. With $\kappa=(64\Phi_*^2)^{-1}$,
\begin{equation}\label{src:eq:gramasym}
 E(n)=\kappa(n-1)^4+O(n^2),\qquad0.068157932354<\kappa<0.068157932355 .
\end{equation}
For all sufficiently large $n$ and every minimizing matrix there is a basis of the lattice generated by its
columns in which, after independent sign changes, the columns are $(1,0)^{\srctr}$ and primitive vectors
$(m,j)^{\srctr}$ on exactly the rows $1\le j\le6$, the number of columns on row $j$ being
\begin{equation}\label{src:eq:gramrows}
 N_j=\frac{w_j}{\Phi_*T_*}\sqrt{1-j^2/T_*^2}\,(n-1)+O(1).
\end{equation}
All implied constants are absolute.
\end{theorem}

The weights in \eqref{src:eq:gramrows} are approximately $0.32484$, $0.15610$, $0.19327$, $0.12773$,
$0.16219$, $0.03587$. B\'ar\'any and Tokushige study the sum of the absolute pairwise determinants in the
problem of the minimum area of a convex lattice polygon \cite{BaranyTokushige2004Minimum}, and Kriepke and
Schymura the largest absolute minor \cite{KriepkeSchymura2025Generic}; for the sum of squares a trace
minimization gives both the scalar optimization and the arithmetic profile.

\emph{Normalization and change of metric.} Let $\Lambda_A=\sum_i\Z v_i$. Its index in $\Z^2$ is $D(A)$, by
the Smith normal form. In a basis of $\Lambda_A$ the columns generate $\Z^2$ and every minor is divided by
$D(A)$ up to a common sign, and Cauchy--Binet gives $\det(AA^{\srctr})=\sum_{i<j}\det(v_i,v_j)^2$. So $E(n)$
is also the least value of the unnormalized sum $\sum_{i<j}\det(v_i,v_j)^2$ over $n$ pairwise nonparallel
integer vectors; the minimum exists, and each column of a minimizing configuration that generates $\Z^2$ is
primitive, since dividing a vector by its content strictly decreases the sum. For a positive definite
$2\times2$ matrix $M$,
\begin{equation}\label{src:eq:gramtrace}
 \min_{g\in\SL_2(\R)}\tr(gMg^{\srctr})=2\sqrt{\det M},
\end{equation}
by the arithmetic--geometric mean inequality, with equality at $g=(\det M)^{1/4}M^{-1/2}$. For a lattice
$\Lambda\subset\R^2$ of covolume one let $\mathcal P(\Lambda)$ contain one vector from each antipodal pair
of primitive vectors. Applying \eqref{src:eq:gramtrace} to $M=\sum_iv_iv_i^{\srctr}$ gives
\begin{equation}\label{src:eq:grammin}
 2\sqrt{E(n)}=\inf_\Lambda\ \min_{S\subset\mathcal P(\Lambda),\ \abs S=n}\ \sum_{v\in S}\abs v^2 .
\end{equation}
Any $n\ge2$ distinct directions span; in a fixed lattice the inner minimum takes $n$ shortest primitive
directions, and an isotropized minimizing matrix realizes the joint infimum.

\emph{Counting primitive directions.} Put
$P_\Lambda(R)=\sum_{v\in\mathcal P(\Lambda)}(R^2-\abs v^2)_+$ and
$N_\Lambda(R)=\#\set{v\in\mathcal P(\Lambda):\abs v\le R}$. After a rotation a reduced basis of $\Lambda$ is
$(a,0),(a\theta,1/a)$, with $a$ the least nonzero length and $\abs\theta\le1/2$; the first vector extends to a
basis, and reducing the second modulo the first gives $a^2\le a^2/4+a^{-2}$, so $a\le(4/3)^{1/4}$. Taking
the representative with positive second coordinate on each nonhorizontal pair, row $j>0$ consists of the
vectors $(a(m+j\theta),j/a)$ with $\gcd(m,j)=1$, and row zero has one representative. Put $T=aR$.

\begin{lemma}[Truncated sums of primitive directions; \stU]\label{src:lem:gramcount}
Uniformly in $\Lambda$, as $R\to\infty$,
\begin{equation}\label{src:eq:gramcrude}
 P_\Lambda(R)=R^4\Phi(T)+O\bigl(R^3\log(2R)\bigr).
\end{equation}
For every compact interval $I\subset(0,\infty)$, uniformly for $T\in I$,
\begin{equation}\label{src:eq:gramsharp}
 P_\Lambda(R)=R^4\Phi(T)+R^2+O_I(1),\qquad N_\Lambda(R)=R^2\beta(T)+1+O_I(1),
\end{equation}
where $\beta(T)=\frac2T\sum_{1\le j\le T}w_j\sqrt{1-j^2/T^2}=2\Phi(T)+\frac T2\Phi'(T)$.
\end{lemma}

\begin{proof}
For $1\le j\le T$ put $\eta=(R^2-j^2/a^2)^{1/2}$. In the horizontal coordinate the summand on row $j$ is
$f(x)=(\eta^2-x^2)_+$, whose integral is $4\eta^3/3$. The mean density of the coprime indices on the row is
$w_j/a$, so the integral terms sum to $R^4\Phi(T)$.

For the uniform error write $\mathbf1_{\gcd(m,j)=1}=\sum_{d\mid\gcd(m,j)}\mu(d)$. On a shifted grid of
spacing $ad$ the sum of $f$ differs from its integral divided by $ad$ by at most the total variation
$2\eta^2$, so the error on row $j$ is $O(R^2\tau(j))$, where $\tau$ counts divisors. Since
$\sum_{j\le T}\tau(j)\le T(1+\log T)$ for $T\ge1$ and $T=O(R)$, summing and adding the horizontal term
proves \eqref{src:eq:gramcrude}, also for $T<1$.

For the finer error, a shifted grid of spacing $d_0$ satisfies
$\abs{\sum_kf(d_0k+u)-d_0^{-1}\int f}\le d_0\operatorname{Var}(f')$: integrate the linear interpolant on
each cell, whose slope is the average of $f'$ there, so the error on a cell is at most $d_0^2$ times the
variation of $f'$ on it. Here $\operatorname{Var}(f')=8\eta$. Split row $j$ into its $\varphi(j)$ primitive
residue classes modulo $j$ and take $d_0=aj$. Only boundedly many rows and classes occur, since $aR=T$ is
bounded, so the total error on positive rows is $O_I(1)$, and the horizontal point contributes
$R^2-a^2=R^2+O_I(1)$. Counting the points of each residue grid in $[-\eta,\eta]$ as its length divided by
$aj$, with error $O(1)$, gives the count. The identity for $\beta$ follows by differentiation, at integers
by continuity.
\end{proof}

\begin{lemma}[The certified scalar maximum; \stU]\label{src:lem:grammax}
The maximum of $\Phi$ is attained at the unique $T_*\in[6.15,6.60]$ with
\begin{equation}\label{src:eq:gramroot}
 \sum_{j=1}^6w_j(4j^2-T_*^2)\sqrt{T_*^2-j^2}=0 ,
\end{equation}
it is nondegenerate, $\Phi(T)<0.4787$ outside $[6.15,6.60]$, and
\[
 6.349265427728<T_*<6.349265427729,\qquad0.478797428690<\Phi_*<0.478797428691 .
\]
\end{lemma}

\begin{proof}
Let $A_0(x)=\sum_{j\le x}w_j=\sum_{d\le x}\mu(d)\floor{x/d}/d$. Comparison with $6x/\pi^2$, bounding the
floor errors by the harmonic sum and the tail by $x/\floor x$, gives $\abs{A_0(x)-6x/\pi^2}\le\log x+3$
for $x\ge1$. Partial summation against $(1-x^2/T^2)^{3/2}$, whose integral over $[0,T]$ is $3\pi T/16$ and
whose total decrease is one, gives
$\abs{\Phi(T)-3/(2\pi)}\le4(\log T+3)/(3T)$ for $T\ge1$, which is below $0.4787-3/(2\pi)$ for $T\ge20000$.
Two finite checks in integer arithmetic close the remaining ranges. First, the totient sieve with upward
rounding of $10^9w_j$ certifies $A_0(j)\le76j/125+3/5$ for $1\le j\le20000$, hence between integers by
monotonicity of the right side, and partial summation gives
$\Phi(T)\le(76/125)\pi/4+4/(5T)<0.4787$ for $1000\le T\le20000$. Second, $[1,1000]\setminus(6.15,6.60)$ is
subdivided into intervals $[l,r]$; the summands of $T\Phi(T)$ are nondecreasing in $T$, so
$\Phi(T)\le r\Phi(r)/l$, and with every radical rounded upwards by integer square roots and bisection of
every interval whose bound is at least $0.4787$, all $3255$ terminal intervals pass. Below $1$ the sum is
empty.

A lower evaluation gives $\Phi(6.35)>0.47879$. With $u=T^2$ the derivative of $\Phi$ on $(6,7)$ has the sign
of $\Gamma(u)=\sum_{j=1}^6w_j(4j^2-u)\sqrt{u-j^2}$, and
$\Gamma''(u)=-\frac{3u}4\sum_{j=1}^6w_j(u-j^2)^{-3/2}<0$. The signs of $\Gamma$ at $u=6.15^2$ and $6.60^2$ are
positive and negative, certified by the same rational bounds on the radicals. Strict concavity gives a
unique zero with $\Gamma'<0$ there, hence $\Phi''(T_*)<0$, and bisection with the monotone bound gives the
intervals for $T_*$, $\Phi_*$ and $\kappa$.
\end{proof}

{\raggedright The finite checks are carried out by
\artifact{gram-euler-and-zero-filters/gram/verify_constant.py}, with output
\artifact{gram-euler-and-zero-filters/gram/constant_verification.json}.\par}

\begin{proof}[Proof of \cref{src:thm:gramasym}]
For every set $S$ of $n$ directions in $\Lambda$ and every $R>0$,
\begin{equation}\label{src:eq:gramduality}
 \sum_{v\in S}\abs v^2\ge nR^2-P_\Lambda(R).
\end{equation}
Take $R^2=(n-1)/(2\Phi_*)$ and $I=[6.15,6.60]$. For $T\in I$, \eqref{src:eq:gramsharp} gives
$P_\Lambda(R)\le\Phi_*R^4+R^2+O(1)$; for $T\notin I$ the error $o(R^4)$ in \eqref{src:eq:gramcrude} is
absorbed, for large $R$, by the gap $\Phi_*-0.4787>0$ of \cref{src:lem:grammax}. With
\eqref{src:eq:grammin},
\[
 2\sqrt{E(n)}\ge(n-1)R^2-\Phi_*R^4-O(1)=\frac{(n-1)^2}{4\Phi_*}-O(1).
\]
For the upper bound take $a=T_*/R$ and the rectangular lattice with basis $(a,0),(0,1/a)$, in which $a$ is
the least length for large $n$. Since $\Phi'(T_*)=0$, \cref{src:lem:gramcount} gives
$N_\Lambda(R)=1+2\Phi_*R^2+O(1)=n+O(1)$. On row one the boundary in the integer coordinate is at
$m_0=(R/a)\sqrt{1-T_*^{-2}}$, and moving $m$ a bounded number of steps past $m_0$ changes
$a^2m^2+a^{-2}$ by $O(1)$, since $aR=T_*$; so the discrepancy $O(1)$ can be corrected with each selected
or discarded point within $O(1)$ of squared radius $R^2$. The resulting set satisfies
$\sum_{v\in S}\abs v^2=nR^2-P_\Lambda(R)+O(1)=(n-1)^2/(4\Phi_*)+O(1)$. Hence
$\sqrt{E(n)}=(n-1)^2/(8\Phi_*)+O(1)$, and squaring gives \eqref{src:eq:gramasym}.

For the structure, isotropize a minimizing configuration by \eqref{src:eq:gramtrace}; its sum of squared
lengths is $2\sqrt{E(n)}$. The gap in \cref{src:lem:grammax} and \eqref{src:eq:gramcrude} force $T\in I$
for large $n$, and the upper bound with \eqref{src:eq:gramduality} and \eqref{src:eq:gramsharp} gives
$R^4(\Phi_*-\Phi(T))=O(1)$, so $T=T_*+O(R^{-2})$ by uniqueness and $\Phi''(T_*)<0$. The nonnegative slack
in \eqref{src:eq:gramduality}, the sum of $R^2-\abs v^2$ over omitted interior points and of $\abs v^2-R^2$
over selected exterior points, is $O(1)$. So every point of squared length below $R^2-O(1)$ is selected and
none above $R^2+O(1)$, with an absolute constant. The horizontal representative is selected, and the rows
$j\ge7$ are excluded since $7/T_*>1$. Each of the rows $1,\dots,6$ meets the circle away from a tangency,
and a bounded change of the squared radius changes its half-width in integer coordinates by $O(1)$,
uniformly in the shear $\theta$. The interval count on row $j$ gives
$N_j=(2w_jR^2/T)\sqrt{1-j^2/T^2}+O(1)$, and replacing $T$ by $T_*$ costs $O(1)$, which gives
\eqref{src:eq:gramrows}; every coefficient there is positive.
\end{proof}

\begin{theorem}[Algebraic degrees of the extremal constants; \stU]\label{src:thm:gramdegree}
The field $K=\Q(T_*^2)$ has degree $48$ and equals $\Q(\kappa)$. Each of $T_*$ and $\Phi_*$ has degree
$96$ over $\Q$.
\end{theorem}

\begin{proof}
Put $u=T_*^2$ and $b_j=30w_j$, so $(b_1,\dots,b_6)=(30,15,20,15,24,10)$. With two parameters $s,t$ let
\[
 B_j=\bigl(b_j(4j^2-x)+s\mathbf1_{j=1}+tb_j(x-j^2)\bigr)\sqrt{x-j^2},
\]
and let $\Omega(x,s,t)$ be the product of $B_1+\sum_{j=2}^6\varepsilon_jB_j$ over the $32$ choices of signs
$\varepsilon_j$. Changing the sign of any radical permutes the factors, up to a common sign change of all
$32$ factors for the first radical; treating the radicals as independent formal square roots, each occurs
with even exponent, so $\Omega\in\Z[x,s,t]$ has degree at most $48$ in $x$. Let $P(x)=\Omega(x,0,0)$ and
let $H_s,H_t$ be the derivatives in $s,t$ at $s=t=0$. On the positive branch at $x=u$ the all-plus factor
vanishes by \eqref{src:eq:gramroot}, so $P(u)=0$.

An exact certificate shows that $\bar P=P\bmod53$ has degree $48$ and is irreducible: $P,H_s,H_t$ are
evaluated at $x=0,\dots,48$ over $\F_{53}$ and interpolated, the product over signs being formed by
multiplying successively by the conjugate under each of the last five sign changes (in the basis indexed by
subsets $I\subseteq\set{1,\dots,6}$ multiplication sends $I,J$ to $I\triangle J$ with the factor
$\prod_{j\in I\cap J}(x-j^2)$, and the product rule carries both derivatives), and in $\F_{53}[X]$
\[
 X^{53^{48}}\equiv X\pmod{\bar P},\qquad\gcd(\bar P,X^{53^d}-X)=1\quad(d=16,24).
\]
As $2,3$ are the prime divisors of $48$, this excludes every proper factor degree, and Gauss's lemma gives
$[K:\Q]=48$.

The root $u$ of $P$ is simple; every radical is analytic near this real point, so exactly one sign factor
vanishes there, and differentiating the product and cancelling the nonvanishing factors gives
\[
 \frac{H_t(u)}{H_s(u)}=\frac{\sum_{j=1}^6b_j(u-j^2)^{3/2}}{\sqrt{u-1}},\qquad
 \kappa=\frac{2025u^4H_s(u)^2}{256(u-1)H_t(u)^2},
\]
with both derivatives nonzero at $u$; so $\kappa\in K$. The leading coefficient of $P$ is a $53$-adic unit,
so reduction maps $\Z_{(53)}[u]$, with the displayed denominator inverted, to $\F_{53}[X]/(\bar P)$, where
the denominator is nonzero; the certificate verifies by cross multiplication that neither the $16$th nor
the $24$th power of Frobenius fixes the image of $\kappa$, which therefore has degree $48$ over $\F_{53}$. A
primitive integer relation for $\kappa$ of smaller degree would reduce to a nonzero relation for this image,
so $\Q(\kappa)=K$.

Let $\ell$ be the leading coefficient of $P$. Since the degree is even, the norms of $u$ and $u-1$ are
$P(0)/\ell$ and $P(1)/\ell$, and the certificate gives residues $50\bmod53$ and $42\bmod59$, quadratic
nonresidues at primes where $\ell$ is a unit. So neither $u$ nor $u-1$ is a square in $K$. Now $T_*=\sqrt u$
and $\Phi_*=2H_t(u)\sqrt{u-1}/(45u^2H_s(u))$; since $\Q(\Phi_*)$ contains $\kappa$ it contains $K$, and both
quadratic extensions are proper.
\end{proof}

The coefficients of $\bar P$ and the degree checks are recorded in the same output file.

For $2\le r\le n-2$ let $E_r(n)$ be the least value of $\det(AA^{\srctr})/D^2$ over integer $r\times n$
matrices $A$ with every maximal minor nonzero, where $D$ is the gcd of their absolute values (the
quantity of \cref{src:thm:gram} with $z=0$ when $n\ge5$), so $E_2(n)=E(n)$.

\begin{corollary}[Contraction and duality; \stU]\label{src:cor:gramrank}
For $2\le r\le n-2$,
\[
 \frac{E_r(n)}{\binom nr}\ge\max\Bigl\{\frac{E(n-r+2)}{\binom{n-r+2}2},\ \frac{E(r+2)}{\binom{r+2}2}\Bigr\},
 \qquad E_{n-2}(n)=E(n).
\]
For fixed $r$ this gives $E_r(n)\ge(2\kappa+o(1))n^2\binom nr$.
\end{corollary}

\begin{proof}
Fix $r-2$ column indices $I$. A unimodular row operation carries the rational span of those columns to the
first $r-2$ coordinate axes, by a basis of its saturated lattice. Every minor extending $I$ is then, up to
sign, a common nonzero integer times a two-row minor of the projected columns; dividing this local family by
its own gcd gives normalized rank-two minors, and the global gcd divides the local one, so the sum of the
squared globally normalized minors extending $I$ is at least $E(n-r+2)$. Each full minor occurs $\binom r2$
times in the sum over $I$; dividing by the number of incidences gives the first bound. The Hodge star of the
primitive exterior product of a basis of the saturated rational row space is, up to sign, the primitive
exterior product for its orthogonal complement, since it spans the right line over $\Q$ and signed
complementary coordinates preserve primitivity. So $E_r(n)=E_{n-r}(n)$, and the first bound applied to
$n-r$ gives the second.
\end{proof}

\begin{example}[Finite witnesses]\label{src:ex:gramwitness}
Rectangular lattices with rational parameters close to $\Phi_*$ and $T_*$, with exact comparison of squared
lengths, give the following upper bounds for $E(n)$; the columns contain $(1,0)^{\srctr}$ and a point on row
one, so their index is one.
\begin{center}\small
\begin{tabular}{rrr}
\toprule
$n$ & $\det(AA^{\srctr})$ & $\det(AA^{\srctr})/(n-1)^4$\\
\midrule
$100$ & $6\,546\,225$ & $0.068147535$\\
$1\,000$ & $67\,885\,760\,907$ & $0.068157984$\\
$10\,000$ & $681\,306\,789\,878\,949$ & $0.068157938$\\
\bottomrule
\end{tabular}
\end{center}
The row counts of the witnesses are recorded with their Gram matrices in
\artifact{gram-euler-and-zero-filters/gram/examples.json}.
\end{example}

\subsection{Why almost all of the target must be spanned}\label{src:sec:span}
\begingroup\raggedright
\sources{grh_lean_transfer_analysis; grh_fixed_field_transfer_2026-11-25; grh_transfer_bridges_2026-11-25; grh_beyond_dirichlet_grh_2026-11-25; grh_beyond_dirichlet_2026-11-25; grh_arithmetic_transfers_2026-11-26; grh_transfer_synthesis_2026-11-25; grh_spanning_and_s5_2026-11-26; grh_transfer_advances_2026-11-26; grh_method_transfers_orthogonal_jets_2026-11-26}
\endgroup

The last line of the source comparison divides by $n$ and needs only $n\ge D/2$. The lower
bound (C) needs much more. In \code{integral_forms_long} each vector completing the selected
family to a basis is taken from a globally spanning auxiliary family and has cost
\[
K_{\mathrm{comp}}=2x_1^2b^2+b ,
\]
the first term bounding its squared norm and the second its denominator height. Before the
rank hypothesis is used, the proof has reached
$0\le c+2\delta+(D-n)K_{\mathrm{comp}}+D(b+u-M)$ with $u=fa\log^4b$, which is the second
hypothesis of \cref{src:prop:rank} with $B=D(b+u)$. Input (B) gives $D-n\le4D/b^2$, and the
lemma \code{mul_completion_cost_le} combines the two:
\[
(D-n)K_{\mathrm{comp}}\le4D\bigl(2x_1^2+b^{-1}\bigr)\le Db
\]
under the scale inequalities of the source. Since $b\le u$, this yields
$c+2\delta\ge D(M-3u)$, which is (C). A fixed positive missing fraction would instead leave a
term of order $Dx_1^2b^2$, far above the error budget. Near-full spanning is thus a
requirement of the lower-bound argument even though the final division uses only a positive
fraction of the dimension, and a statement that a fixed fraction of the target is spanned
cannot be substituted for (B) unless the completion is made correspondingly cheap. The
quantitative spanning statements, including a sharper form of (B) and the weakest completion
budget under which the source argument still runs, are in \cref{sec:spn}.

\begin{example}[Determinants do not control individual directions]\label{src:ex:diag}
The positive forms $\operatorname{diag}(t,t^{-1})$ and the identity on $\R^2$ have the same
determinant, while their restrictions to the first coordinate line have ratio $t$, which is
arbitrary. A bound for a full determinant therefore controls neither the determinant of a
subfamily nor an individual eigenvalue. This is why the lower bound must retain the completing
vectors and their costs, and why the spanning theorem, which controls a large family through
a second moment, does not isolate a prescribed eigenform.
\end{example}

The length $m\asymp\sqrt x$ of the chain is tied to the denominator budget in the same way.
The budget allotted to local defects in the fitting of twists is $fa(\log b)^3/(mk)$, with
$a=xb$ and a further parameter $k$ of the construction, and the explicit estimate for these
defects is $O\bigl(f(m/k)\,b\log b\bigr)$. The quotient of the estimate by the budget is
\begin{equation}\label{src:eq:budget-ratio}
O\Bigl(\frac{m^2}{x(\log b)^2}\Bigr),
\end{equation}
so $m\asymp\sqrt x$ leaves only logarithmic room. A longer chain with the same scale choices
would need better loss estimates, and whenever a normalised loss grows like $m$ the division
by $m$ leaves a constant and the decay disappears. The exponent in \eqref{src:eq:onesided}
cannot be improved by enlarging $m$ with the remaining estimates held fixed.

\subsection{Quantifiers in the twisting step}\label{src:sec:quantifiers}
\begingroup\raggedright
\sources{grh_transfer_analysis; grh_transfer_analysis_2026-11-25; grh_beyond_dirichlet_grh_2026-11-25; grh_beyond_dirichlet_2026-11-25; grh_fixed_field_transfer_2026-11-25; grh_method_transfers_2026-11-25; grh_method_transfers_2026-11-25_1; grh_method_transfers_2026-11-25_2; grh_method_transfers_orthogonal_jets_2026-11-26; grh_transfer_bridges_2026-11-25; grh_transfer_extensions_2026-11-26; grh_ideal_densities_and_transfer_conditions_2026-11-26; grh_arithmetic_transfers_2026-11-26; grh_structural_transfers_2026-11-25}
\endgroup

The lower bound \eqref{src:eq:twisted} and the twisting theorem behind it are stated with
quantifiers that are easily blurred. The following distinctions are part of the mechanism.

\paragraph{Local twists, one global height.}
In \eqref{src:eq:twisted} the favourable twist may depend on the rational prime and on the
embedding into an algebraic closure of $\Q_p$; no single twist is required to be good at every
finite place. What links the local choices is the predicate \code{HeightLe}, which contains a
\emph{common} nonzero multiplier together with bounds at all archimedean conjugates and at
all finite places, and this is what makes the product formula applicable to the ratios
$Z(1)/Z(\xi)$. In addition, the matrix that clears the index at a given $p$-adic
identification must work for all twists allowed there, and integrality is required at an
embedding and at its conjugate. Independently favourable local objects without this common
height control do not define a global lattice inequality. Conversely, the hypothesis of
\eqref{src:eq:twisted} that every $\Theta_\xi$ is injective is used in its proof only at the
trivial twist and at the twists actually selected, where it makes the reference Gram
determinant and the selected determinant ratios nonzero; injectivity at the selected twists
suffices.

\paragraph{Products and individual valuations.}
A bound for a product of determinants bounds the sum of their valuations, not each valuation:
an unexpectedly small valuation of one factor can offset an inadmissibly large one of another.
The selection results of \cref{sec:tors} therefore keep the individual inequalities and
bound the bad choices by a union bound. The same care is needed when finite certificates are
aggregated over several probes, since a witness may change from probe to probe.

\begin{lemma}[Products before gcds; \stU]\label{src:lem:gcd}
Let $\ell$ be a prime and let $(A_{q,j})$ be a finite array of integers indexed by a nonempty
finite set of probes $q$ and a nonempty finite set of branches $j$. Then
\[
\bigl(\forall q\ \exists j:\ \ell\mid A_{q,j}\bigr)
\quad\Longleftrightarrow\quad
\ell\ \Bigm|\ \gcd_q\Bigl(\prod_jA_{q,j}\Bigr) .
\]
\end{lemma}

\begin{proof}
For fixed $q$, primality of $\ell$ shows that $\ell$ divides the product over $j$ exactly when
it divides some factor. Dividing every product is equivalent to dividing their gcd.
\end{proof}

\begin{example}\label{src:ex:gcd}
For two probes and two branches with rows $(A_{1,1},A_{2,1})=(2,3)$ and
$(A_{1,2},A_{2,2})=(3,2)$, each branch has gcd $1$ over the probes, but the products over
branches are $(6,6)$: the primes $2$ and $3$ survive every probe by switching branch. Eliminating
each branch by a different probe is therefore insufficient unless a theorem providing a single
common witness has first been proved.
\end{example}

\paragraph{The conclusion of the twisting theorem.}
For a fixed coefficient prime $\ell$, finitely many Selmer data and a power range
$1\le e\le m$, the theorem \code{exists_anticyclotomic_twist} produces a finite set of
auxiliary primes and one finite-order anticyclotomic character $\psi$ whose first $m$ powers
simultaneously have bounded normalised Selmer lengths, controlled valuations of the attached
value systems and small auxiliary ramification cost. In the readable edition the predicate
\code{TwistedLengthLe} bounds the length of the Selmer module with divisible coefficients
$\operatorname{Frac}(O')/O'$, for all sufficiently large coefficient rings $O'$. The
intermediate torsion precision is chosen with a strict margin above the error bound, and a
torsion module whose $\pi^N$-torsion has length less than $N$ is equal to that torsion
(\code{length_selmerModule_divisibleQuot_le}); so the conclusion is not restricted to a finite
level and no interchange of ``for every precision there is a twist'' with ``there is a twist
for every precision'' occurs. The details, and the comparison with the finite-precision
statement displayed in the compact edition, are in \cref{sec:sel}. The remaining
qualifications are arithmetic.
\begin{itemize}
\item A \code{SelmerDatum} requires residual absolute irreducibility, a local rank budget,
finite conjugate-dual and local defects, and an inertial element acting by a scalar $\zeta$
with $\zeta^2-1$ a unit; a nonzero pure Tate twist, being unramified away from $\ell$, cannot
satisfy the last condition directly (\cref{sec:sel}). The local conditions of a datum are
those of the source and must be compared with the conditions of any arithmetic group one has
in mind.
\item The order of $\psi$, the coefficient extension and the auxiliary primes may vary with
the requested tolerance. The theorem bounds an order exponent and normalised losses. It does
not bound the norms of the auxiliary primes or the conductor of $\psi$, and it does not
prescribe a fixed order or give one character good at every coefficient prime.
\item The auxiliary characters are of the form $\psi=\xi/\xi^c$ with $c$ complex conjugation,
where the precursor $\xi$ is ramified on only one side of each conjugate pair and is trivial
on the whole decomposition groups, not only the inertia groups, at the prescribed places
(\code{IsFromU1}). Triviality of $\psi$ alone at those places is a weaker condition. The
effect on which local torsion data extend to global twists, and on their conductors, is worked
out in \cref{sec:spn}.
\item The value systems enter through a seed whose valuation is assumed finite; the theorem
preserves nonvanishing of an existing seed and does not create it.
\item Lengths are normalised by the absolute ramification index of the coefficient ring. For
$O_N=\Z_\ell[\zeta_{\ell^N}]$ with uniformiser $\varpi_N=\zeta_{\ell^N}-1$, the nonzero module
$O_N/(\varpi_N)$ has normalised length $1/(\ell^{N-1}(\ell-1))$, which tends to $0$. A
normalised bound that shrinks over increasingly ramified coefficient rings therefore does
not give vanishing; over a fixed coefficient ring a length less than $1$ does.
\end{itemize}

\subsection{What a transfer must supply, and what does not transfer}\label{src:sec:limits}
\begingroup\raggedright
\sources{grh_transfer_analysis; grh_transfer_analysis_2026-11-25; grh_lean_transfer_analysis; grh_beyond_dirichlet_grh_2026-11-25; grh_fixed_field_transfer_2026-11-25; grh_method_transfers_2026-11-25; grh_method_transfers_2026-11-25_2; grh_transfer_synthesis_2026-11-25; grh_transfer_advances_2026-11-26; grh_arithmetic_transfers_2026-11-26; grh_ideal_densities_and_transfer_conditions_2026-11-26; grh_structural_transfers_2026-11-25; grh_a6_transfer_2026-11-25; grh_cyclic_generator_transfer_2026-11-25; grh_deficit_transfer_2026-11-25; grh_relative_prime_races_2026-11-25; grh_sylow_transfer_2026-11-25; psl2_grh_transfer; grh-lean-transfers}
\endgroup

Three kinds of reuse occur in this article. A \emph{consequence} uses \cref{src:thm:dgrh} as a
black box (\cref{sec:zer,sec:dir,sec:cls,sec:smd}). A \emph{transfer of an ingredient}
develops one component of the proof for another problem and usually does not depend on
\cref{src:thm:dgrh} at all (\cref{sec:tors,sec:cost,sec:lub,sec:gra,sec:gpp,sec:gpi,sec:spn}).
A \emph{transfer of the mechanism} would re-create the comparison of
\cref{src:sec:mechanism} for a new prime statistic, and is the route to zero-free half-planes
for $L$-functions outside \cref{src:thm:dgrh}. For the last kind the analytic step is the
portable part: \cref{sec:ana,sec:nab,sec:grp,sec:hec} give one-sided criteria and detecting
weights for many targets. The arithmetic part is not available for any weight other than
$\lambda_{d,q}$. What it would have to consist of can be stated exactly.

\begin{proposition}[Sufficient arithmetic package; \stR]\label{src:prop:package}
Let $R_w$ and $R_*$ be real functions on $[1,\infty)$, and let $A,C,C'>0$ and $\eta>0$.
Suppose that for some $x\ge1$ there are real numbers $m\ge\sqrt x$ and $0<n\le d$, positive
reals $a,\tau,f$ and $b>1$ with $\log b\le x^{\eta}$, and reals $h,V,\delta,\lambda$ and
$x_1\ge1$ such that, with
\[
\srcB=fa\log^3b,\qquad L=\srcB\log b,\qquad
\srcRR=a\tau\bigl(R_w(x)-R_w(x_1)+\lambda R_*(x_1)\bigr),\qquad h_{\mathrm{top}}=h-m\srcRR+L,
\]
the following hold:
\begin{gather*}
(1-4/b^2)d\le n,\qquad d(h-3L)\le V+2\delta,\qquad V\le n\,h_{\mathrm{top}},\qquad
\delta\le5d\log(1+b^2)\srcB,\\
-h_{\mathrm{top}}\le b^2x^{A},\qquad 4x^{A}+4L+10\log(1+b^2)\srcB\le C'(\log b)^2\srcB,\qquad
f/\tau\le C\log b,\\
R_w(x_1)-\lambda R_*(x_1)\le x^{-1}.
\end{gather*}
Then
\[
R_w(x)\le CC'\,\frac{(\log b)^6}{m}+x^{-1}\le CC'x^{-1/2+6\eta}+x^{-1} .
\]
\end{proposition}

\begin{proof}
\Cref{src:prop:top} with $\kappa=4/b^2$, $T=b^2x^A$ and $J=5\log(1+b^2)\srcB$ gives
$m\srcRR\le4L+4x^A+10\log(1+b^2)\srcB\le C'(\log b)^2\srcB=C'fa(\log b)^5$. Dividing by
$ma\tau>0$ and using $f/\tau\le C\log b$,
\[
R_w(x)-R_w(x_1)+\lambda R_*(x_1)\le CC'\frac{(\log b)^6}{m} .
\]
Add the last hypothesis and use $m\ge\sqrt x$ and $\log b\le x^\eta$.
\end{proof}

For the weight $\lambda_{d,q_0}$ the source constructs such data for every small $\eta$ and all
large $x$, with $R_w=R_{d,q_0}$, $R_*=R_{d,1}$, $\lambda=\tau_0/\tau$ and $A=99$. If the
hypotheses were established, for every small $\eta$ and all large $x$, for the centred ramp
sum $R_w$ of a weight $w$ detecting the zeros of a target $L$-function, the criteria of
\cref{sec:ana} would give the corresponding zero-free half-plane $\Re s>1/2$; if they were
established only for one fixed $\eta$, the conclusion would reach only the corresponding
line. No such data are constructed in this article for any $w$ other than the source weight,
and \cref{src:prop:package} is a reduction, not a proof of a Riemann hypothesis. The
second-scale hypothesis $R_w(x_1)-\lambda R_*(x_1)\le x^{-1}$ can be discharged by a window
theorem (\cref{ana:thm:window}) when the comparison is proved for every $x_1$ in a window and
the coefficient $\lambda$ is fixed independently of $x_1$. Two cautions apply at this point:
upper bounds for a remainder at two scales cannot be subtracted to bound their difference, and
small averages of several errors do not give a point at which all of them are small (for
$E_1=\sin$ and $E_2=\cos$ one has $E_1^2+E_2^2=1$).

A construction meeting the hypotheses of \cref{src:prop:package} for a new weight has to
provide all of the following, for the same lattice.
\begin{enumerate}[label=\textup{(G\arabic*)}]
\item \emph{A prime detector realised by divisibility.} Integral forms whose local
divisibility records the required Frobenius or Hecke data, at all good primes and with control
of ramified primes and prime powers, together with an archimedean norm formula having the
same main term. In the source the occurrence of $\Lambda(n)(1-\chi_d(n))\mathbf 1_{n\equiv1(q)}$
is proved through the local divisibility of theta derivatives and matching global periods; a
class function of a nonabelian Galois group, or a Hecke trace, is not a parameter of that
theorem. For a weight with coefficients of both signs, the comparison must be arranged so
that the cancellation takes place between the upper and the lower bound, since a negative
gain in valuation is not available for integral sections.
\item \emph{Short forms and enough of them.} The analogues of (A) and (B) for the new field,
level, lattice and weights, with a rank deficit and a completion cost compatible with
\cref{src:prop:rank}.
\item \emph{A compatible lower bound.} The analogue of (C) for the same integral structure and
chain of lifts, with the same main term, the determinant heights of
\cref{src:sec:quantifiers} and a uniform total denominator budget; the hypotheses of the
twisting theorem must be verified for the new local types.
\item \emph{Scales.} Existence of admissible data throughout the range of the second cutoff,
with $m\asymp\sqrt x$ and total normalised error $x^{o(1)}$.
\item \emph{The analytic endgame.} A transform whose poles are the zeros in question, with
the sign of real residues under control. When the residues are differences of multiplicities,
as in \eqref{src:eq:residue}, a device preventing cancellation is needed; when the statistic
comes from a product all of whose zero multiplicities have the same sign, that step
disappears. The detectors of \cref{sec:nab,sec:grp} are of the second kind, which removes the
auxiliary zero-density input but none of (G1)--(G4).
\end{enumerate}

The following limits apply to every such attempt.

\paragraph{Congruence weights see only the abelian quotient.}
The weights of the source depend on a rational residue class and a quadratic character. A
finite combination of Dirichlet characters factors through an abelian quotient of the Galois
group, and for a Galois extension $E/\Q$ with nonabelian group $G$ every element of the
commutator subgroup acts trivially on $E\cap\Q(\zeta_M)$ for every $M$; by the Chebotarev
density theorem a class function taking different values at the identity and at some element
of the commutator subgroup cannot agree, on all large unramified primes, with a function of
$p\bmod M$. The quantitative form of this obstruction is in \cref{sec:nab}. Consequently (G1)
needs a genuinely new carrier of Frobenius information; retuning the modulus or the
cyclotomic data in the identity for $\srcrem$ cannot supply it.

\paragraph{Fixed targets and uniformity.}
The estimates \eqref{src:eq:onesided-int} and \eqref{src:eq:onesided} are eventual estimates
for a fixed modulus and a fixed auxiliary prime, with no uniformity in the conductor. Problems
that need uniformity over growing moduli or over families of fields, such as the
Elliott--Halberstam conjecture or Hooley's treatment of Artin's primitive-root conjecture
through the Kummer fields $\Q(\zeta_k,a^{1/k})$, are not reached by a fixed-target estimate.
For the first, summing the bound
$\psi(x;q,a)-x/\varphi(q)=O(\sqrt x\log^2(qx))$ of Dirichlet GRH over $q\le Q$ gives only
$O(Q\sqrt x\log^2(xQ))$. For the second, the Kummer fields are in general nonabelian over $\Q$;
see \cref{sec:grp}.

\paragraph{Location of zeros is not multiplicity or correlation.}
A function can have all its zeros on a line and a zero of any multiplicity there, as
$(s-\tfrac12)^k$ shows. \Cref{src:thm:dgrh} and the mechanism behind it give no information on
simplicity of zeros, on vanishing or nonvanishing at the central point, or on correlations
between zeros. The prime sums controlled are sums in one variable; correlations such as those
of $\Lambda(n)\Lambda(n+2)$, needed for twin primes or binary Goldbach, are not addressed.

\paragraph{Determinants and second moments are averages.}
By \cref{src:ex:diag} a determinant bound over a large space does not bound an individual
form. Applications to individual automorphic $L$-values, to subconvexity or to
Ramanujan-type bounds would need an isolation step that the source does not contain.

\paragraph{Internal notions are specialised.}
Several background theories inside the source are developed only in the strength its proof
needs: the attachment of Galois representations is a statement about $p$-adic limits along a
special filter of Frobenius data, compatibility at a split place says that a Hecke polynomial
annihilates Frobenius and does not by itself identify the characteristic polynomial, and no
general Weil--Deligne construction is supplied. A result about these internal objects is a
statement about a standard arithmetic object only after a comparison has been proved;
\cref{sec:sel} gives such comparisons for the odd cocycle module and records what remains
for Selmer groups, the Birch--Swinnerton-Dyer conjecture, Vandiver's conjecture and algebraic
$K$-groups. In particular, the shortened source bounds the lengths of odd cocycle modules for
a fixed cyclotomic CM field by the valuations of one integer, and its denominator estimate is
arranged so that the even classes do not enter; avoiding those classes is not a proof that
they vanish. The remark of the original announcement on $K$-groups and a weakening of
Vandiver's conjecture \cite[Step~17]{Alpoge2026Remark} refers to an interface of the original,
larger development that is absent from the shortened source tree.

\subsection{Open questions}\label{src:sec:open}
\begingroup\raggedright
\sources{grh-lean-transfers; grh_method_transfers_2026-11-25; grh_fixed_field_transfer_2026-11-25; grh_transfer_bridges_2026-11-25; grh_relative_prime_races_2026-11-25; grh_a6_transfer_2026-11-25; grh_beyond_dirichlet_grh_2026-11-25; grh_new_results_2026-11-26; gram-euler-and-zero-filters}
\endgroup

\begin{question}\label{src:q:divisibility}
Is there a family of integral forms whose local divisibility at an unramified prime $p$
distinguishes Frobenius classes that have the same image in every abelian quotient, for
instance the identity and the $3$-cycles of the splitting field of $X^3-X-1$ inside the same
quadratic splitting condition, with an archimedean norm formula exhibiting the same main term?
This is requirement (G1) in the smallest nonabelian case.
\end{question}

\begin{question}\label{src:q:package}
For one fixed nonabelian target, can the hypotheses of \cref{src:prop:package} be established
with all four costs recorded: divisibility gain, rank deficit, completion cost and denominator
loss?
\end{question}

\begin{question}\label{src:q:completion}
In a new setting, can integral completing vectors be constructed whose excess cost $K-A$ in
\cref{src:prop:rank} is of the size of the other errors, so that a fixed positive fraction of
the target suffices in place of near-full spanning?
\end{question}

\begin{question}\label{src:q:constantfree}
Can a comparison of covolumes be set up whose remainder is directly a difference of the prime
statistic at two nearby scales, so that neither the centring constant nor a second, remote
cutoff is needed?
\end{question}

\begin{question}\label{src:q:gram}
For fixed $(r,n,z)$, what is the least value of $\det(AA^{\srctr})/D^2$ in
\cref{src:thm:gram}? The three affine bounds are each attained for $r=2$, $n=5$, but their
maximum is not claimed to be optimal for every triple. For $r=2$ and $z=0$,
\cref{src:thm:gramasym} gives $E(n)=\kappa(n-1)^4+O(n^2)$: what are the limit points of
$(E(n)-\kappa(n-1)^4)/n^2$, at which the row endpoints and the basis shear enter, and from which $n$
on are all minimizers confined to six rows? For $r=2$ and a fixed number of vanishing minors, where
the multiplicities of repeated directions must also be optimized, what is the asymptotic minimum? For
fixed $r\ge3$, what is the sharp constant, or even the order of growth, of $E_r(n)/\binom nr$, which
\cref{src:cor:gramrank} bounds below by $(2\kappa+o(1))n^2$?
\end{question}
\endgroup
